% Chapter 2. Exact original source carriers; evidence class: SOURCE DEFINITION.
\chapter{Carriers, Indices, Function Spaces, and Notation}
\label{ch:carriers}

\section{Why the Hilbert space must be specified first}
The equations of matrix quantum mechanics can be written compactly, but their
mathematical content depends upon the spaces on which those equations act. A
fermionic matrix expression, a smooth wavefunction and an equivalence class of
square-integrable wavefunctions are distinct kinds of object. Many later
questions---unitary implementation, charge closure, the existence of adjoints,
and spectral confinement---are meaningful only after these distinctions are
fixed. We therefore use the \emph{literal source carriers} from
\path{OAI.MathematicalPhysics.BFSS.Core}, not a replacement representation
chosen for computational convenience.

Let $N$ denote the size of the color matrices, and set
\begin{equation}\label{eq:source-indices}
 \mathsf{Space}=\{0,\ldots,8\},\qquad
 \mathsf{Spin}=\{0,\ldots,15\},\qquad
 \mathsf{Color}_{N}=\{0,\ldots,N^2-2\}.
\end{equation}
Here $\mathsf{Color}_{N}$ has $N^2-1$ elements. The research program of the
present volume specializes to $N=2$, so there are exactly three color labels.
The nine spatial labels specify a nine-dimensional vector representation;
they are \emph{not} a claim that our quantum mechanics has eleven independent
coordinates or that it is equivalent to eleven-dimensional M-theory.

\begin{definition}[Original source bosonic carrier]\label{def:boson-carrier}
For $N\geq2$, the source's real Euclidean configuration space is
\begin{equation}\label{eq:boson-carrier}
 \mathcal X_N
 :=\mathbb R^{\mathsf{Space}\times\mathsf{Color}_N}
 \cong\mathbb R^{9(N^2-1)}.
\end{equation}
Its coordinates are $x_i^A$, with $i\in\mathsf{Space}$ and
$A\in\mathsf{Color}_N$.
\end{definition}
The distinguished case is
\begin{equation}\label{eq:boson-27}
 \mathcal X:=\mathcal X_2\cong\mathbb R^{27},
 \qquad x=(x_i^A)_{i=0,\ldots,8}^{A=0,1,2}.
\end{equation}
Thus 27 counts bosonic \emph{matrix components}, not spacetime dimensions.
The potential's commuting-matrix valleys are geometrically special subsets
of this configuration space; their existence will later prevent us from
assuming that a quartic bosonic potential is globally radially confining.

\begin{definition}[Original source fermionic carrier]\label{def:fermion-carrier}
The complex Hilbert fiber is
\begin{equation}\label{eq:fermion-carrier}
 \mathcal F_N:=\mathbb C^{\,2^{8(N^2-1)}}.
\end{equation}
For $N=2$ this is $\mathcal F:=\mathcal F_2=\mathbb C^{2^{24}}$.
The source uses the ordinary Euclidean Hermitian Hilbert structure and
represents each fermionic generator by a bounded complex-linear operator on
this \emph{same} $\mathcal F$.
\end{definition}
The interpretation of the exponent is important: sixteen real generators
for each of three colors give 48 real Majorana generators, which are paired
into 24 complex creation--annihilation modes. The finite-dimensional complex
Fock representation has dimension $2^{24}$, not 48 and not $2^{48}$.
We verify that operator-algebra identification in Chapter~\ref{ch:clifford};
for now the exact fiber is part of the source definition.

\section[The L2 Hilbert space is a different carrier]{The $L^2$ Hilbert space is a different carrier}
\begin{definition}[Original full wavefunction Hilbert space]
Equip $\mathcal X_N$ with Lebesgue volume $dx$. The original complex Hilbert
space of square-integrable Fock-valued wavefunctions is
\begin{equation}\label{eq:full-l2}
 \mathscr H_N:=L^2(\mathcal X_N,dx;\mathcal F_N),\qquad
 \mathscr H:=\mathscr H_2=L^2(\mathbb R^{27};\mathbb C^{2^{24}}).
\end{equation}
A member is an \emph{almost-everywhere equivalence class}, not a particular
pointwise function. A convenient representative is written $x\mapsto\psi(x)$
only when the distinction matters.
\end{definition}

There is no finite-volume hypothesis here:
$\operatorname{vol}(\mathbb R^{27})=\infty$. Thus a uniformly bounded
operator field $x\mapsto B(x)\in\mathcal L_{\mathbb C}(\mathcal F)$ need not
itself be square integrable. In contrast, the product $B(x)\psi(x)$ is
square integrable when $\psi\in\mathscr H$ and
$\sup_x\|B(x)\|<\infty$, provided the product is strongly measurable.
The distinction underlies the separate normalized-field multiplier program
of Chapter~\ref{ch:norm}; it cannot be bypassed by pretending that all
bounded functions belong to $L^2$.

\begin{proposition}[Essentially bounded fields on the original $L^2$ carrier;
conventional functional analysis]\label{prop:bounded-l2-multiplier}
Let $B:\mathcal X\to\operatorname{End}_{\mathbb C}(\mathcal F)$ be a
strongly measurable field such that
$\operatorname*{ess\,sup}_{x\in\mathcal X}\|B(x)\|\leq C<\infty$.
Then pointwise multiplication
$(M_B\psi)(x)=B(x)\psi(x)$, interpreted almost everywhere,
defines a bounded complex-linear operator
$M_B:\mathscr H\to\mathscr H$ with $\|M_B\|\leq C$.
\end{proposition}
\begin{proof}
In the finite-dimensional fiber, the product of a measurable operator
field and a measurable vector representative is measurable. If two
representatives of $\psi$ agree almost everywhere, their products
with $B$ do likewise. The map is therefore well-defined on $L^2$
equivalence classes, and
\[
 \|M_B\psi\|_{\mathscr H}^{2}
 =\int_{\mathcal X}\|B(x)\psi(x)\|_{\mathcal F}^{2}\,dx
 \leq C^2\int_{\mathcal X}\|\psi(x)\|_{\mathcal F}^{2}\,dx.
\]
Linearity passes from pointwise linearity to the equivalence classes.
\end{proof}
This is a conventional functional-analytic statement, not a new source
Lean certificate. The measurability and essential bound must be established
for an \emph{actual source field} before it can be applied there. An
auxiliary endomorphism-valued construction, or a bound at a single
configuration, does not supply those hypotheses on the original
Fock-\emph{vector} Hilbert space.

\begin{definition}[Original smooth test-function core]
The original source core is
\begin{equation}\label{eq:smooth-core}
 \mathscr D_N:=C_c^\infty(\mathcal X_N;\mathcal F_N).
\end{equation}
The canonical map $\iota:\mathscr D_N\rightarrow\mathscr H_N$ sends a
compactly supported smooth function to its $L^2$ equivalence class. In the
source it is named \texttt{coreToL2}.
\end{definition}

\begin{proposition}[Core functions define actual $L^2$ vectors]
For each $\psi\in\mathscr D_N$, the $L^2$ class $\iota\psi$ is defined,
and addition and complex scalar multiplication commute with $\iota$.
\end{proposition}
\begin{proof}
Continuity and compact support imply square integrability of a
finite-dimensional vector-valued function on Euclidean space. Addition and
scalar multiplication preserve its almost-everywhere equivalence class, so
$\iota(\psi+\phi)=\iota\psi+\iota\phi$ and
$\iota(c\psi)=c\iota\psi$. The source implements these identities using
\texttt{MemLp.toLp\_add} and \texttt{MemLp.toLp\_const\_smul}; this is a
proof about the chosen original Hilbert carrier, not a new physical model.
\end{proof}
\formalrecord{Definitions \texttt{Boson}, \texttt{Fermion}, \texttt{FullL2},
\texttt{SmoothCore}, and \texttt{AlgebraData.coreToL2} in
\texttt{Core.lean}, source snapshot \texttt{etblink/math@d4870eb1ac58}.
Atlas records D-01--D-03.}

\section{Conventional inner products, coordinate operators, and a warning}
The original configuration coordinate $x_i^A$ acts on a smooth core
wavefunction by multiplication; its directional derivative is
$\partial_{iA}$. These are \emph{unbounded} operators in the full
$\mathscr H$, even though they are well-defined as linear operations on
$\mathscr D$. In particular, the fact that a first-order differential
expression maps test functions to test functions does not make its closure,
adjoint, or squared Hamiltonian automatic.

For later use we write
\begin{equation}\label{eq:original-inner}
 \langle \psi,\phi\rangle_{\mathscr H}
 :=\int_{\mathcal X}\langle \psi(x),\phi(x)\rangle_{\mathcal F}\,dx,
 \quad \|\psi\|_{\mathscr H}^2
 =\int_{\mathcal X}\|\psi(x)\|_{\mathcal F}^2\,dx.
\end{equation}
All weak-form equations later in this volume follow the source's
\texttt{inner} conventions. We will not infer the domain of an unbounded
operator from an identity on $\mathscr D$, nor infer equality of $L^2$
classes from equality at a single point.

\section{Reader's checkpoint}
At this stage we have introduced only the setting: nine spatial labels,
three SU(2) color labels, the 48-real-generator/24-complex-mode fermion
fiber, the true vector-valued $L^2$ Hilbert space, and its compactly
supported smooth core. None of these choices establishes an eigenstate,
a spatial symmetry on the physical space, or an M-theory interpretation.
The next task is to explain what the fermionic operators are and why the
particular $2^{24}$-dimensional representation is the source one.
